% Solutions to the Chapter 7 exercises that are not code, from lectures/L07/appendix/c_exercises.md.
% Exercises 7.6 and 7.7 are Code and are not answered; see chapter.tex for why.

\section{Chapter 7: Concurrency and Locking}
\label{sol:sec:c7}
Answers to the exercises in \secref{c7:sec:exercises}. Exercises~\ref{c7:ex:6} and~\ref{c7:ex:7}
are Code, and are checked on the target as their Check yourself lines say.

% ------------------------------------------------------------------------------------------
\solution[fillaccent2]{c7:ex:1}{Where concurrency comes from}{Recall}

\xpart{a} SMP, preemption, interrupts, timers and workqueues (\secref{c7:app:a1}). On a single-CPU
kernel built with \code{PREEMPT_NONE}, two remain: interrupts, and timers, because a
\code{timer_list} callback runs in softirq context on the way out of an interrupt and so can run in
the middle of your code just as a handler can. SMP goes with the second CPU and preemption with
\code{PREEMPT_NONE}. A workqueue item runs in a kernel thread, which on such a kernel gets the CPU
only when your code gives it up by sleeping or returning, so it can overlap you only across a sleep.

\xpart{b} That nothing else runs this code at the same time, and that it cannot be stopped
mid-function. So it protects shared data, if at all, by disabling interrupts on its own CPU, which
does nothing about another CPU running the same code at the same moment, and its read-modify-write
sequences can be preempted halfway with another task running the same code in between. What breaks
is data shared between callers, silently and intermittently: lost updates, a ring buffer whose level
no longer matches its contents. The fix is a lock, which the old code has no notion of.

\xpart{c} Interrupts. The answer does not change on a single CPU: the handler can run on top of
\code{read()} on the same CPU at any point where interrupts are enabled, which is the case
\code{spin_lock_irqsave} exists for (\secref{c7:app:a5}). A second CPU adds SMP as well, since the
handler can then also run on the other CPU while \code{read()} runs, and that is what the lock
itself excludes. \code{spin_lock_irqsave} covers both.

\xpart{d} Three: a load of the counter into a register, an add, and a store back. The add happens in
a register, where no other CPU can see it. What matters is the other CPU's load and store falling
between this CPU's load and its store: if both load $n$ before either stores, both store $n+1$, and
one of the two increments is lost (\secref{c7:app:a2}).

% ------------------------------------------------------------------------------------------
\solution[fillaccent2]{c7:ex:2}{Atomics and their limits}{Recall}

\xpart{a} The ring buffer of Chapter~\ref{c5:ch}. Its invariant: \code{level} is between 0 and the
capacity, and the \code{level} bytes starting at \code{head}, wrapping at the capacity, are exactly
the bytes written and not yet read, in the order they were written. It mentions the head, the level
and the array at once, so making each of them atomic leaves it breakable between any two of their
updates (\secref{c7:app:a3}); Exercise~\ref{c5:ex:8} constructs an interleaving that does it.

\xpart{b} Make \code{head} and \code{level} plain integers and put the whole test, store and
increment inside one critical section: take a spinlock; if \code{level < capacity}, store the byte
at \code{(head + level) \% capacity} and increment \code{level}; release the lock. A spinlock,
because the section is a handful of instructions that never sleep and so can be held briefly from
any context; \code{spin_lock_irqsave} if an interrupt handler also reads or writes the buffer. A
mutex would be correct too if every caller is in process context, but buys nothing for a section
this short. The point is that the test and the update become one indivisible step, so no second
writer can pass the same test and pick the same slot in between.

\xpart{c} Atomics protect one variable; locks protect an invariant: if what must be true before and
after mentions more than one variable, you need a lock.

\xpart{d} The compiler may load \code{flag} once, before the loop. Nothing in the body can change
\code{flag} as far as the compiler can see, no store, no call and no barrier, so reading it again
would be redundant, and the loop spins on a register forever. The fix is
\code{while (!READ_ONCE(flag)) cpu_relax();}, with the thread that sets it using
\code{WRITE_ONCE(flag, 1)}. It is not a memory barrier: \code{READ_ONCE} forces one real load on
every iteration and orders nothing on the CPU. If the waiter must then read data the other thread
wrote before setting the flag, that needs ordering as well: \code{smp_store_release} to set it and
\code{smp_load_acquire} to read it, or a lock.

% ------------------------------------------------------------------------------------------
\solution[fillaccent3]{c7:ex:3}{Choosing a primitive}{Hand calculation}

\xpart{a} A spinlock: two kernel threads run in process context, no handler touches the counter,
and the section is one increment, far too short to be worth sleeping over. (A mutex would also be
correct, and an \code{atomic_t} better still, but that is not one of the choices.)

\xpart{b} A spinlock with \code{_irqsave}: the handler takes the lock in interrupt context, so
\code{read()} must keep interrupts off its own CPU while it holds it, and \code{irqsave} does so
without assuming the state it was called in.

\xpart{c} A mutex: \code{copy_to_user} may fault and sleep, so the section may sleep, and only a
mutex may be held across a sleep.

\xpart{d} A spinlock: 2 microseconds is short, nothing in the sequence sleeps, and one sequence must
not be interleaved with another CPU's. If the interrupt handler also touches those registers, it is
\code{_irqsave}, as in b).

\xpart{e} A spinlock with \code{_irqsave}: the timer callback runs in softirq context, so
\code{write()} must keep it off its own CPU while it holds the lock, and with interrupts off no
softirq runs there either. (The exact minimum is \code{spin_lock_bh} in \code{write()} and a plain
\code{spin_lock} in the callback.)

\xpart{f} A mutex: an allocation that may have to reclaim is \code{GFP_KERNEL} and may sleep, so
the lock held across it must be one a sleeper may hold. The alternative is to allocate before
taking a spinlock.

\xpart{g} b) and e). In b), \code{read()} holds the plain spinlock, the device interrupts that CPU,
and the handler spins on the lock; the holder cannot run again until the handler returns, and the
handler never returns, so that CPU is gone, and so is every other CPU that then wants the lock. In
e) it is the same with the timer: the timer interrupt arrives while \code{write()} holds the lock,
the timer softirq runs the callback on the way out of it, on the same CPU, and the callback spins on
a lock whose holder is underneath it. (A plain spinlock in c) or f) is a different bug, sleeping
while holding one, which \code{CONFIG_DEBUG_ATOMIC_SLEEP} reports.)

% ------------------------------------------------------------------------------------------
\solution[fillaccent3]{c7:ex:4}{What the race costs}{Hand calculation}

\xpart{a} $4 \times 20{,}000 = 80{,}000$.

\xpart{b} $80{,}000 - 63{,}562 = 16{,}438$ lost, and $16{,}438 / 80{,}000 = 20.5\%$.

\xpart{c} The range is $30{,}047 - 1{,}564 = 28{,}483$ increments, and
$28{,}483 / 80{,}000 = 35.6\%$ of the expected total: the runs lost from $1{,}564 / 80{,}000 =
2.0\%$ to $30{,}047 / 80{,}000 = 37.6\%$. The spread says that the loss is not a property of the
program but of the timing of one run. A single run is one draw from a very wide distribution and
predicts nothing about the next: the same code on the same machine lost 2\% once and 38\% another
time, so one run can neither size the bug nor show whether a change made it better.

\xpart{d} Zero: yes. A lost update needs one thread's store to land between another thread's load
and store; if the threads never run their loops at the same time, that never happens and the count
is exact. It has been seen: the first version of the module, without the barrier, lost nothing
three runs in a row, and at 200 and 2,000 increments per thread every run was exact
(\secref{c7:app:a2}).

60,000: yes, in principle, although nothing close to it was observed. A lost update is a stale
store: a thread writes back the value it loaded plus one, wiping out every increment made since its
load, not just one. A thread preempted between its load and its store, on a \code{CONFIG_PREEMPT}
target running four threads on two CPUs, can come back after the others have added thousands and
erase all of them with a single store. The final value can in fact be as low as 2: thread A loads 0;
the others run everything except B's last increment, reaching 59,999; A stores 1; B loads 1; A runs
its remaining 19,999, reaching 20,000; B stores 2. So a loss of 60,000, a final count of 20,000, is
possible. The largest loss measured was 30,047.

\xpart{e} First, one sample of a nondeterministic result. A single exact run is what the broken
program produces some of the time, and at 2,000 per thread it produced it every time, so it says
nothing about the next run. Second, taking a run as the argument for correctness. \code{counter++}
on a shared variable from several threads is three unprotected steps, and that is a race by
construction whatever any run returns: whether the lock is needed follows from the code, and an
observation can show a bug but never its absence.

% ------------------------------------------------------------------------------------------
\solution[fillaccent4]{c7:ex:5}{Atomic context}{Design}

\xpart{a} Not atomic: \code{read()} runs in process context for the calling task. Legal.

\xpart{b} Atomic, in an interrupt handler. Illegal: \code{GFP_KERNEL} may sleep to reclaim.

\xpart{c} Atomic. Legal: \code{GFP_ATOMIC} never sleeps, and the handler must cope with
\code{NULL}.

\xpart{d} Atomic, since a spinlock is held. Illegal: \code{mutex_lock} sleeps when the mutex is
contended.

\xpart{e} Not atomic: holding a mutex does not make you atomic, and you are still in process context.
Legal; the section inside the spinlock is then atomic, and nothing in it may sleep.

\xpart{f} Atomic. Illegal: \code{copy_to_user} may fault, and bringing the page in may sleep.

\xpart{g} Atomic: a timer callback runs in softirq context. Illegal: \code{msleep} sleeps, every
time.

\xpart{h} b), d) and f). g) never works: \code{msleep} always calls into the scheduler, and from
softirq context that prints ``BUG: scheduling while atomic'' on the first run. The three that
usually work:
\begin{itemize}
  \item \textbf{b)} works because the allocation is served from a slab cache without reclaim. It
    fails the day memory is short and the allocator has to reclaim, which means sleeping inside a
    handler.
  \item \textbf{d)} works because an uncontended \code{mutex_lock} is a single atomic
    compare-and-exchange and never sleeps. It fails the first time another task holds the mutex,
    and the spinlock holder goes to sleep.
  \item \textbf{f)} works because the user's buffer is resident, so the copy never faults. It fails
    when the page is not there, never touched or since reclaimed, and the copy has to fault it in,
    which may sleep. On a kernel that counts preemption, as this course's does, the fault handler
    sees that it is in atomic context and declines the fault (\code{faulthandler_disabled()},
    \file{include/linux/uaccess.h}), so the copy comes up short and the driver returns
    \code{-EFAULT} instead. Either way it happens only under memory pressure.
\end{itemize}

\xpart{i} \code{CONFIG_DEBUG_ATOMIC_SLEEP}. It puts a check at the top of every call that might
sleep, before the call decides whether it has to: \code{might_alloc()} in \code{kmalloc},
\code{might_sleep()} in \code{mutex_lock}, and \code{might_fault()} in \code{copy_to_user}. All
three print ``BUG: sleeping function called from invalid context'' the first time they run in
atomic context, on a machine with plenty of memory, an uncontended mutex and a resident page.

% ------------------------------------------------------------------------------------------
\solution[fillaccent4]{c7:ex:8}{What the tools do not see}{Design}

\xpart{a} Lockdep watches lock acquisitions. At each one it records which locks were already held,
adds the resulting edges to its graph and checks them, and checks the context the lock is taken in
(\secref{c7:app:b3}). \code{qa_race} with \code{locked=0} takes no lock at all, so as far as lockdep
is concerned nothing happens: no acquisition, no edge, no context to check. The lost updates are
plain loads and stores to a shared variable, of which lockdep has no model; it cannot know that the
variable needed a lock.

\xpart{b} Lockdep checks the locks you take, never the locks you did not take. A codebase that
``passes with lockdep enabled'' has shown that on the paths its tests ran, the locks it took were
taken in a consistent order and in legal contexts. That says nothing about data that should have
had a lock and had none, and nothing about paths the tests never ran.

\xpart{c}
\begin{itemize}
  \item \textbf{Waiting for an event rather than a lock:} two threads each wait for a completion,
    or a wait queue condition, that only the other will signal. The graph is built from lock
    acquisitions, and waiting for an event is not one, so the cycle has no edges.
  \item \textbf{Waiting on hardware:} a driver holds a lock while it polls for a status bit, or
    waits for an interrupt, that never comes. The other party in the cycle is a device, which
    lockdep cannot see at all.
\end{itemize}

\xpart{d} The error paths: an allocation that fails, a device that times out or answers wrongly or
disappears, a probe that fails halfway and unwinds, a signal that interrupts a sleep, a remove while
a file is still open. A test suite drives the normal path because that is what it can cause;
failures have to be injected, and mostly are not. It matters more for lock ordering because an
inversion takes two orders and lockdep reports only once it has seen both, and the unusual order
typically lives on exactly those paths: cleanup that takes locks in the reverse of the order set-up
took them. The failure is also not confined to the rare path. Most bugs on an unexercised path hurt
only the call that takes it; an ordering bug deadlocks the machine when the rare path meets the
common one.

\xpart{e} The module's licence is not GPL-compatible, so loading it tainted the kernel, and that
taint is one that switches lock debugging off for the rest of the boot: the module loader adds
\code{TAINT_PROPRIETARY_MODULE} with \code{LOCKDEP_NOW_UNRELIABLE}, and \code{add_taint} then turns
lockdep off (\file{kernel/module/main.c}, \file{kernel/panic.c}). The line is
\code{Disabling lock debugging due to kernel taint}. You first met it in Chapter~\ref{c4:ch}, in
Exercise~\ref{c4:ex:7}.

% ------------------------------------------------------------------------------------------
\solution[fillaccent]{c7:ex:9}{A number computed twice}{Cross-check}

\xpart{a} $4 \times 20{,}000 = 80{,}000$.

\xpart{b} A correct module prints \code{actual=80000 lost=0} on every run; with the spinlock
switched on, the course's runs all returned exactly 80,000 (\secref{c7:app:a2}).

\xpart{d} Worked here with the course's ten unlocked runs from \secref{c7:app:a2} in place of your
five.
\begin{itemize}
  \item With the lock, the result is determined by the program rather than sampled from it: each of
    the 80,000 increments is a whole load, add and store that no other increment can interleave
    with, so the counter must end at $4 \times 20{,}000$, and 79,999 would be a bug, not noise.
    Nothing is measured on a scale that has an error; the program counts. Chapter~\ref{c3:ch}'s
    \code{Image} size is an exact count too, but not one you can predict: it depends on the
    configuration, the compiler and the source, so you measure it, and a different build
    legitimately gives a different number. Chapter~\ref{c2:ch}'s context-switch rate is a rate
    over a sampling interval, which varies with the load and with the interval you happen to catch,
    so two correct measurements differ. Here the prediction is arithmetic, and the measurement must
    agree with it to the last digit.
  \item The ten losses sum to 154,616, so the mean loss is 15,461.6, which is 19.3\%, and the mean
    \code{actual} is 64,538.4. The range is 1,564 to 30,047, a spread of 28,483, or 35.6\% of the
    expected total. The mean on its own is not useful: the spread is nearly twice the mean
    ($28{,}483 / 15{,}461.6 = 1.84$), and both depend on the conditions, the number of threads and
    CPUs, the increments per thread, the barrier and the load. To mean anything it needs the spread,
    the number of runs and the conditions beside it; and what matters for the bug is not the average
    loss but that the loss is not zero.
  \item If every collision lost exactly one increment, the worst run, 30,047 lost, took 30,047
    collisions. But a thread preempted between its load and its store erases every increment made
    while it was off the CPU, so a single collision can lose thousands. Each collision loses at
    least one, so 30,047 is an upper bound on the number of collisions, not a count. Even as a
    count it is plausible: 30,047 in a few hundred milliseconds is of the order of 100,000 a
    second ($30{,}047 / 0.3\,\mathrm{s} \approx 100{,}000$), for two CPUs doing nothing but load,
    add and store on one shared variable. The collision rate is limited by the rate of increments,
    not by the clock.
\end{itemize}

\xpart{e} The course's twelve runs (\secref{c7:app:a2}): at 200 and at 2,000 per thread, 0\%, 0\%
and 0\%; at 20,000, 23.3\%, 20.8\% and 6.9\%; at 200,000, 31.7\%, 34.9\% and 26.5\%.
\begin{itemize}
  \item The loss appears at 20,000 and at 200,000. At 200 and 2,000 every run is exact.
  \item The bug is not absent at that size; the code is identically broken. What differs is that at
    2,000 per thread no two threads are in their loops at the same time for long enough to collide:
    each thread's work is over before the others have got going, so the race has no window to
    happen in.
  \item Released together means released from the barrier together, not running together. After the
    release each thread still has to be woken and scheduled onto a CPU, and with four threads on two
    CPUs two of them wait for a CPU in any case. The overlap has to last long enough for two threads
    to be executing their loops on the two CPUs at once, or for one to be preempted between its load
    and its store, and it is competing against the wake-up and scheduling latency after the release:
    at 200 or 2,000 increments the whole loop is shorter than that stagger. The measurements put the
    threshold between 2,000 and 20,000 increments per thread, and \secref{c7:app:a2} estimates it
    at about ten thousand. How long that is in time depends on this machine's increment rate and its
    scheduling latency, which the chapter does not publish; time a run near the threshold to find
    out.
  \item They would conclude that the code is correct, and they would be very confident, every run
    exact, run after run. They would be wrong, and the confidence would come from repeating a test
    that could not exercise the race, which is no evidence at all.
\end{itemize}

\xpart{f} The bug report gets the unpredictable one: the unlocked losses, with their spread and
the conditions they came from (four threads, 20,000 increments each, two CPUs, the barrier),
because a report's job is to show that the bug exists, how bad it is, and how to reproduce it. The
test gets the predictable one: \code{locked=1} must give exactly 80,000, because a test has to pass
or fail the same way on the same code every time, and only an exact prediction can be asserted
exactly. An assertion on the unlocked run can only be ``lost more than nothing'', and even that is
safe only because the barrier and 20,000 increments made a zero loss something never measured,
which is why \file{run.sh} runs it three times and asks only that at least one run lose something.

\xpart{g} For this program, yes. The only shared state is one counter, and \code{atomic_inc} makes
its load, add and store indivisible, so the count comes out at exactly 80,000. With a running
maximum alongside, no. The counter and a maximum that must agree with it form an invariant over two
variables, and updating the maximum is a read, compare and write separate from the increment: two
threads can each read, compare and store in the wrong order, and a reader can see one variable
updated and not the other. Per-variable atomicity does not make the sequence atomic.
\Secref{c7:app:a3} answers it: atomics protect one variable, locks protect an invariant.
